\addchap{Vorwort und Hinweise}
\label{ch:Vorwort}

\subsection*{Nutzung externer Tools}
Für die sprachliche Ausarbeitung und Korrektur der Einleitung (\cch{Einleitung}) und der Durchführung (\cch{Durchfuehrung}) wurden die KI \textit{Lumo AI} sowie der \textit{Duden Mentor} eingesetzt. Die inhaltliche Auswertung und die eigentliche Erstellung des Dokuments erfolgten jedoch vollständig ohne Unterstützung durch Künstliche Intelligenz. \cite{Duden,Lumo}

Alle im Dokument enthaltenen Grafiken ohne Referenz wurden eigenständig in Inkscape erstellt bzw. nachbearbeitet oder sind in Python geplottete Grafiken. Der Python-Code sowie das LaTeX-Dokument wurden in der Entwicklungsumgebung Visual Studio Code (VS Code) angefertigt.

% Im Bereich der Code-Erstellung wurde KI unterstützend eingesetzt. Sollte eine Code-Zelle primär durch eine KI generiert worden sein, so wurde dies explizit innerhalb der jeweiligen Zelle vermerkt, um Transparenz zu gewährleisten.

\subsection*{Eigener Code und Transparenz}
Zur Durchführung der Auswertungen wurde eine eigene Python-Bibliothek entwickelt. Der gesamte Quellcode - sowohl die Python-Skripte als auch die LaTeX-Quellen - ist selbst erstellt und öffentlich online einsehbar.

Die gesamte Projekt-Struktur ist auf im GitHub zu finden. Hier ist sowohl der Python-Code, als auch der Latex-Code zu finden. \cite{githubPAP2}

Die benutzten Pythonfunktionen sind gesonders im GitHub unter dem Namen\afz{Papulator} zu finden. \cite{Papulator}

\subsection*{Verwendung universitärer Materialien}
Dieses Dokument basiert auf dem Skript von Jens Wagner (Stand: 11. Februar 2026). Für die Einleitung wurden Grafiken aus diesem Skript übernommen. Diese Entnahmen sind transparent gekennzeichnet; die verwendeten Abbildungen Stammen immer aus den Kapiteln des Skriptes,welche dieselbe Versuchsnummer und denselben Versuchsnamen wie im vorliegenden Dokument tragen. In diesem Versuch wurden die Grafiken aus \emph{Versuch \versuchsnummer \versuchsname{\versuchsnummer}} verwendet. \cite{skriptPAP21}